\section{复杂度与课程算法思想}

NPQR 的复杂度主要来自候选映射、候选 SWAP 展开、束搜索和局部修复。令 $n=|V|$
表示物理量子比特数，$m$ 表示电路门数，$e=|E|$ 表示硬件边数，$K$ 表示初始映射
候选数，$B$ 表示束宽，$D$ 表示主搜索最大展开深度，$A$ 表示每个状态保留的候选
动作数，$R$ 表示后缀修复深度。

\begin{table}[H]
  \centering
  \caption{复杂度}
  \label{tab:complexity}
  \begin{tabularx}{\linewidth}{lYY}
    \toprule
    阶段 & 时间复杂度估计 & 说明 \\
    \midrule
    拓扑距离矩阵 & $O(n^3)$ 或 $O(n(e+n\log n))$ & 取决于 Floyd 或多源最短路；固定拓扑下可复用。\\
    线路依赖图构建 & $O(m)$ & 扫描门序列，维护前沿门和执行依赖。\\
    逻辑交互图构建 & $O(m)$ & 统计双比特门交互强度，服务初始映射。\\
    初始映射候选 & $O(K\cdot n^2)$ 量级 & 与候选数量和映射评分方式有关。\\
    动作生成与掩码 & $O(B\cdot e)$ & 每个保留状态在硬件边上筛选候选 SWAP。\\
    模型推理 & $O(B\cdot A\cdot C_\theta)$ & $C_\theta$ 为单次模型评分成本。\\
    束搜索扩展 & $O(D\cdot B\cdot A\cdot C_s)$ & $C_s$ 包括状态复制、执行门和评分更新。\\
    后缀修复 & $O(A^R)$ 的受限搜索 & 只在接近完成时启动，避免全局穷举。\\
    轨迹复放 & $O(m+N_{\mathrm{swap}})$ & 线性验证输出轨迹是否合法。\\
    \bottomrule
  \end{tabularx}
\end{table}


复杂度表明，若不加剪枝，搜索空间会随深度指数增长。NPQR 通过固定束宽、动作
剪枝和有界修复把搜索规模限制在可控范围内。因此，该算法属于近似搜索框架：它
追求有限时间内的高质量可行解，而不保证全局最优。

\subsection{课程算法思想映射}

项目可以对应到多种课程算法设计思想。表 \ref{tab:course-mapping} 给出主要对应
关系。这种映射不是把每个经典算法完整复刻一遍，而是说明 NPQR 在问题建模和工程
实现中如何使用这些思想。

\input{tables/course_mapping.tex}

\subsection{SABRE 的基线意义}

SABRE 的核心思想可以看作贪婪启发式与有限前瞻的结合。它根据前沿门和后续门在
物理拓扑上的距离选择 SWAP，通常能以较低时间成本给出稳定路线。NPQR 与 SABRE 的
区别在于：NPQR 引入神经网络动作评分，并在搜索中保留多条候选路线，同时加入局部
修复和轨迹复放。两者可以用于对比，但不能把 SABRE 结果作为 NPQR 的成功结果。
